\documentclass[10pt,a4paper]{article}
\usepackage{fontspec,amsmath,amssymb,graphicx}
\usepackage[margin=18mm]{geometry}
\setmainfont{Times New Roman}
\setlength{\parindent}{0pt}
\setlength{\parskip}{5pt}
\allowdisplaybreaks
\sloppy
\emergencystretch=2em
\begin{document}
\section*{BT\_5\_5\_3\_03 — Lời giải bổ sung}
Hàm là nhân Poisson với cực $(1,0)$. Nó điều hòa trong đĩa mở, dương, nhưng không liên tục và không bị chặn ở cực: $u(r,0)=(1+r)/(1-r)\to\infty$. Ở các điểm biên khác cực, giới hạn bằng 0. Nguyên lý cực đại đòi liên tục trên miền đóng nên không áp dụng như đề phát biểu cho quả cầu đóng chứa cực. Không có mâu thuẫn: hàm không đạt cực đại/cực tiểu trong miền mở.\end{document}